\documentclass[11pt]{article}
\usepackage{amsmath,amssymb}
\usepackage[margin=1in]{geometry}
\usepackage[hidelinks]{hyperref}
\title{Counting full covers in Erd\H{o}s problem 687(i)}
\author{Research notes for the Left to Solve submission}
\date{November 2026}
\begin{document}
\maketitle

For $P=\{p\text{ prime}:p\le x\}$, write $a=(a_p)_{p\in P}$ for a choice of
one residue modulo each prime. Given a finite set $S$ of integers, let
$U(a)$ be the set of its elements missed by all the chosen classes. For
$S=[1,H]$, let $N_x(H)$ count the assignments with $U(a)=\varnothing$.
The maximum in the problem satisfies $H\le Y(x)$ exactly when $N_x(H)>0$.
Thus a candidate scale $F(x)$ can be established by proving positivity at
$H=\lfloor cF(x)\rfloor$ and vanishing at $H=\lceil CF(x)\rceil$ for fixed
$c,C>0$ and every sufficiently large $x$, provided $F(x)\to\infty$.

\section*{An exact reduction to residue counts}
Write $\nu_p(T)=|\{n\bmod p:n\in T\}|$. For any integer $t$,
\[
 \sum_a(1+t)^{|U(a)|}
 =\sum_{T\subseteq S}t^{|T|}\prod_{p\in P}\bigl(p-\nu_p(T)\bigr).
\]
Indeed the summand on the left expands as
$\sum_{T\subseteq U(a)}t^{|T|}$. For prescribed $T$, the number of
assignments missing all its elements is exactly the displayed product,
since the residues at different primes are chosen independently. Taking
$t=-1$ and $S=[1,H]$ gives an integer formula for $N_x(H)$. Restricting
the same double count to $|T|=k$ gives every factorial moment,
\[
 \sum_a \binom{|U(a)|}{k}
 =\sum_{\substack{T\subseteq S\\ |T|=k}}\prod_{p\in P}\bigl(p-\nu_p(T)\bigr).
\]
The alternating sum at $t=-1$ must be controlled with its sign.
The following sufficient upper test instead uses nonnegative terms. For
positive integers $m,r$, if
\[
 \sum_a \bigl||U(a)|-m\bigr|^{2r}<m^{2r},
\]
then no assignment covers $S$, since a covering assignment contributes
$m^{2r}$ by itself. The sums here are over all assignments, without
probability normalization. These finite identities leave open the estimates
needed for all large $x$.

The connection between moments of reduced residues and the Jacobsthal
function has an established literature. See Montgomery and Vaughan,
\emph{On the distribution of reduced residues} (1986),
\url{https://annals.math.princeton.edu/1986/123-2/p04}, and Arnaud
Chadozeau, \emph{Sur la répartition des entiers premiers à un entier
donné} (2006), especially the discussion of bounds uniform in the order
of the moment, \url{https://oskar-bordeaux.fr/handle/20.500.12278/25339}.

For the order-of-growth question the upper test can be weakened further.
If an assignment covers $[1,y]$ with $y\ge H$, each of its $y-H+1$
subintervals of length $H$ translates to a cover of $[1,H]$. The resulting
assignments are distinct. Two equal shifts would differ by a multiple
of every prime up to $x$, hence by a multiple of their product, whereas
their positive difference is at most $y$, strictly below that product.
Consequently $Y(x)<H+N_x(H)$. In particular $N_x(H)\le H$ suffices for
$Y(x)<2H$. The moment test for this conclusion is
$\sum_a\bigl||U(a)|-m\bigr|^{2r}\le Hm^{2r}$, since each full cover
contributes exactly $m^{2r}$ to this nonnegative sum.
Define $T(x)$ as the first $H$ with $N_x(H)\le H$. Since
$N_x(Y(x)+1)=0$, we have $T(x)\le Y(x)+1$, and the preceding bound
gives $Y(x)<2T(x)$. Thus an estimate for the growth of this counting
threshold would give the order of $Y(x)$ itself. The finite formulas do
not yet determine the growth of $T(x)$.

\section*{What the second moment already gives}
Take the uniform distribution on all assignments for a finite set of
primes $P$. Put $A=\prod_{p\in P}p$, $V=\prod_{p\in P}(1-1/p)$, and
let $U$ be the number of missed positions among $0,\ldots,H-1$.
Then $\mathbb E U=HV$. For positions $i,j$, the number of assignments
missing both is
\[
 \sum_{\substack{D\subseteq P\\d_D\mid i-j}}
      \prod_{p\in P\setminus D}(p-2),\qquad d_D=\prod_{p\in D}p.
\]
This follows by expanding $\prod_{p\in P}(p-2+\mathbf1_{p\mid i-j})$.
For each fixed $i$ there are at most $H/d_D+1$ choices of $j$
in its class. Summing over $i$ and multiplying out the two products
gives $\mathbb E U^2\le (HV)^2+HV$. Thus
\[
       \operatorname{Var}(U)\le HV,\qquad
       \mathbb P(U=0)\le 1/(HV)\quad (H>0).
\]
The translation $i\mapsto i+1$ permutes all assignments and gives the
same distribution on $[1,H]$. This classical second-moment estimate
controls the proportion of full covers. The absolute count still has
the factor $A$. The finite divisor expansion and the variance and count
bounds are in \texttt{Variance687.lean}.

\section*{Small-prime choices admitting an adaptive completion}
Fix $m\ge144$, $z=k^m$, $x=k^{m+1}$ and $H=z^2$. Choose independently
one uniform residue for each $p\le z$. Let $E$ be the event that
\emph{some} assignment at the primes $z<p\le x$ completes the cover.
Write $X$ for the number of survivors at $z$ and
$V(t)=\prod_{p\le t}(1-1/p)$. For sufficiently large $k$,
the square-window deletion estimate gives $X\le(24/m)H/\log z$ on $E$.
As $V(z)\ge1/(3\log z)$, this is at most half the mean $HV(z)$.
The variance estimate above implies
\[
   \Pr(E)\le \frac{4}{HV(z)}\le\frac{12\log z}{z^2}.
\]
The same estimate holds on $[0,H)$ by translating all residues.
Lean 4.33.1 checks the uniform second-moment bound on this event in
\texttt{ProfilsProlongeables687.rare}.

A deterministic normalization also applies to an optimal cover. For a
set $P$ of prescribed prime rows and arbitrary residues $b_p$, put
$D=\prod_{p\in P}p$. The Chinese remainder theorem supplies $0\le t<D$
with $t\equiv a_p-b_p\bmod p$ for every $p\in P$ in an optimal assignment
$a$. Translate every row by $-t$ and keep the interval
$[1,Y(x)-(D-1)]$, interpreted as empty if the upper endpoint is too
small. All its positions remain covered and every prescribed row now
agrees with $b_p$. In particular, for $x\ge1$ one can prescribe
\emph{all} primes through $\lfloor\log_{16}x\rfloor$ and retain a cover
of length $Y(x)-\lfloor\sqrt x\rfloor$, since $D\le4^{\lfloor\log_{16}x\rfloor}$.
Lean verifies \texttt{ClassesImposees687.logarithmic\_prefix} for every
prescription. This permits, for example, testing only zero residues at
these primes when an additive cost of $\lfloor\sqrt x\rfloor$ is acceptable.

An analytic argument gives a stronger conditional estimate. Prescribe
arbitrarily every class through $w=\lfloor\log_{16}z\rfloor$ and choose
independent uniform classes on $(w,z]$. Write $D=\prod_{p\le w}p$.
Since $D\le4^w\le\sqrt z$, complete periods of length $D$ yield at least
$(H-D)V(w)$ survivors of the prescribed classes. Each of these survives
the random choices with probability $V(z)/V(w)$. Thus the conditional
mean of $X$ exceeds $H/(6\log z)$ eventually. For $m\ge288$, the event
$E$ entails a downward deviation of at least $H/(12\log z)$.
By the square-window stability estimate, changing just the class at
$p\in(w,z]$ affects $X$ by at most $68H/(p\log z)$.
The \href{https://doi.org/10.1017/CBO9781107359949.008}{bounded differences
inequality of McDiarmid} gives, uniformly in the prescribed prefix,
\[
  \Pr(E\mid (a_p)_{p\le w})
  \le \exp\!\left(-\frac{1}{72\cdot68^2
                       \sum_{w<p\le z}1/p^2}\right).
\]
Summing dyadically using $\pi(2t)\le16t/\log t$ gives
$\sum_{p>w}p^{-2}\le32/(w\log w)$ for $w\ge2$. The bound is therefore
$\exp(-c\log z\log\log z)$ for some absolute $c>0$ and all sufficiently
large $k$. Using its finite exponential-moment bound, Lean 4.33.1
also verifies, for every $w\ge2$ with $2w\le z$ and
$\prod_{p\le w}p\le z$, the explicit estimate
\[
  \Pr(E\mid (a_p)_{p\le w})
      \le \exp\!\left(-\frac{w\log w}{2\cdot12^2\cdot68^2\cdot32}\right).
\]
The declarations are \texttt{ConcentrationUniforme687.uniform} and
\texttt{ConcentrationUniforme687.logarithmic}. The latter instantiates
$w=\lfloor\log_{16}z\rfloor$. The number of
possible assignments to the middle primes still exceeds the reciprocal
of the displayed bound, so exceptional profiles can persist.

The same calculation permits deeper conditioning. For fixed $0<a<1$,
put $w=\lfloor z^a\rfloor$. The Chinese remainder theorem converts any
prescribed prefix into a translated interval sieved by the zero classes.
The lower linear sieve at parameter $2/a>2$ therefore gives at least
$c_a H V(w)$ survivors, uniformly in the prefix, for a constant $c_a>0$.
One may use a fixed sieve parameter strictly between $2$ and $2/a$ to
absorb the integer cutoff. See the interval bound in
\href{https://arxiv.org/abs/2010.01211}{Granville, formula (1)}.
The conditional mean is at least $c_a H V(z)$. Taking $m\ge144/c_a$
puts the threshold for completion below half that mean. The same
bounded differences argument gives
\[
    \Pr(E\mid (a_p)_{p\le w})\le \exp(-c'_a z^a\log z)
\]
for all sufficiently large $k$, with $c'_a>0$. Here $a$ and $m$ must
be fixed before $k$ tends to infinity.

\section*{A necessary distribution for nearly quadratic covers}
Suppose there is a sequence $x\to\infty$ of covers of length $H$ with
$\log H/\log x\to2$. Fix $1/2<\alpha<1$ and put $z=\lfloor x^\alpha\rfloor$.
Write $S$ for the positions avoiding the chosen classes through $z$,
$N_p$ for the number of elements of $S$ hit by the chosen class at $p$,
and $m(n)$ for the number of hits at $z<p\le x$. Then
\[
 \frac{|S|\log x}{H}\longrightarrow \log\frac{2-\alpha}{\alpha},
 \qquad
 \frac{\log x}{H}\sum_{n\in S}(m(n)-1)\longrightarrow0 .
\]
In particular the positions whose first covering prime lies between
$x^\alpha$ and $x^\beta$, for fixed $\alpha<\beta<1$, have limiting
count, in units of $H/\log x$,
$\log((2-\alpha)/\alpha)-\log((2-\beta)/\beta)$.
The limits apply to every such sequence of assignments.

Indeed set $s=\log H/\log z$ and $r=\log x/\log z$.
For $2<s<4$, the lower interval sieve gives, in units $H/\log z$,
the lower envelope $L=(2/s)\log(s-1)+o(1)$ for $|S|$.
The indices in one progression modulo $p$ form an interval of length
$H/p+O(1)$. Multiplication by $p$ permutes the classes of each
prime up to $z$, so that CRT converts their restrictions to ordinary
interval sieving. At $t=\log p/\log z$ the upper interval sieve
therefore gives $N_p\le U_p$, where
$U_p/(H/\log z)=(2/(s-t)+o(1))/p$. Here $s-t$ stays in a compact subset of
$(1,3)$ for each fixed $\alpha$ in the indicated range.
These constants follow from formula (1) and the explicit sieve
functions in \href{https://arxiv.org/pdf/2010.01211}{Granville,
\emph{Sieving intervals and Siegel zeros}, pp.\ 1--2}. Uniformity
over a compact range of exponents also follows by dividing the
range into finitely many subintervals, varying the sieve cutoff
monotonically at fixed interval length, then taking the mesh to zero.

Partial summation using Mertens' theorem gives the normalized sum of
row envelopes as $\int_1^r 2\,dt/(t(s-t))+o(1)$. Subtracting the
lower envelope yields exactly
\[
 \int_1^r \frac{2\,dt}{t(s-t)}
       -\frac2s\log(s-1)=\frac2s\log\frac{r}{s-r}.
\]
The right side tends to zero since $s/r\to2$. The excess of the
actual stock over the lower envelope, the sum of the row deficits
$\sum_p(U_p-N_p)$, and the number of repeated hits are all
nonnegative after including the sieve errors in the envelopes.
Their sum is the total envelope difference, by counting all hits in
two orders. Each normalized quantity therefore tends to zero.
Multiplying $L$ by $\log x/\log z$ gives the claimed stock limit.
Lean verifies the finite accounting identity, the squeeze, and the
exact integral as \texttt{RigiditeCritique687.balance},
\texttt{.saturation}, and \texttt{.sieve\_gap}.
In particular a positive normalized excess in any one of these
quantities, uniformly along putative nearly quadratic covers, would
exclude that sequence. No such uniform excess is established here.

The conclusions also hold if $o(H/\log x)$ positions are left
uncovered. In that case count repeated hits as
$\sum_{n\in S}\max(m(n)-1,0)$. Each uncovered position has $m(n)=0$
and adds one unit to the available envelope difference, which
changes its normalized value by $o(1)$.
This tolerance matters for a proof strategy. Under the hypothesis of
infinitely many Siegel zeros, Granville's Corollary 1 at the fixed
parameter $2$ gives a sequence of intervals of length $H=x^{2+o(1)}$
with $o(H/\log x)$ survivors after sieving through $x$.
Consequently the global stock law and the vanishing normalized
repeat count can hold for these nearly complete assignments too.
An argument that excludes every such numerical profile would also
exclude that exceptional-zero hypothesis.

For complete covers there is a local form. Let $Q=Q(x)\ge1$ satisfy
$\log Q=o(\log x)$. For a residue $b$ modulo $Q$ avoiding $a_p$
at every $p\mid Q$, the same calculation gives
\[
 \#\{n\in S:n\equiv b\bmod Q\}
   =\left(\log\frac{2-\alpha}{\alpha}+o(1)\right)
       \frac{H}{\varphi(Q)\log x}.
\]
If $b$ hits one of those prescribed classes the count is zero. On an
allowed progression the new interval length is $H/Q+O(1)$, with
the same limiting logarithmic exponent $2$. The excluded primes
dividing $Q$ are all below $z$ eventually. Their deletion from
the sieve multiplies its Euler product by $Q/\varphi(Q)$.
This same product occurs in the upper envelopes after restricting
to a larger prime, since each such prime is coprime to $Q$.
Iwaniec's upper and lower corollaries for excluded divisors of $Q$
have relative additive bracket error
$O(\log\log(3Q)/\log N)=o(1)$ at each interval length $N$ used here.
Their constants are uniform in $Q$ and in the interval origin.
See the \href{https://github.com/jasondavies/iwaniec-linear-sieve/blob/84dfeff69393f31bbdcd126531f0ad8976069dad/IwaniecLinearSieve/Main/Corollary.lean}{formal statements of both corollaries}.
Thus the global squeeze applies inside every such progression, even
when the modulus and the admissible class vary along the sequence.

\paragraph{The least prime at every fixed power cutoff.}
The stock law extends to every fixed $0<\alpha<1$. Let $f,F$ be
the usual dimension-one sieve functions in Granville's formula (1),
with $(u f(u))'=F(u-1)$ for $u>2$ and $f(2)=0$. Under the same
near-quadratic hypothesis one has
\[
 \frac{\log x}{H}|S_\alpha|
      \longrightarrow C(\alpha):=\frac{e^{-\gamma}}{\alpha}f(2/\alpha).
\]
Here $S_\alpha$ consists of positions with no hit at $p\le x^\alpha$.
Assign each remaining covered position to its \emph{least} hitting
prime $p$. Within its progression the smaller rows sieve an interval
of length $H/p+O(1)$ through $p^-$, the cutoff just below $p$.
Writing $t=\log p/\log x$, the upper sieve and Mertens' product give
the sum of the normalized row envelopes as
\[
 e^{-\gamma}\int_\alpha^1 F(2/t-1)\frac{dt}{t^2}
   =\frac{e^{-\gamma}}2\int_2^{2/\alpha} F(u-1)\,du
   =C(\alpha).
\]
The lower sieve for $S_\alpha$ gives exactly this constant.
The functions are continuous on the compact ranges used; replacing
$H/p$ by its integer part has uniformly vanishing relative error.
If $H/p<p$, sieving the row only through its length gives the same
limiting upper envelope at $t=1$, where $F(1)=2e^\gamma$.
Thus the normalized sum of the deficits of the least-prime rows
tends to zero. When finitely many fixed cutoffs are used these
conclusions hold simultaneously. The argument allows $o(H/\log x)$
uncovered positions, which contribute directly to the counting identity.

There is a genuine change in the number of repeated hits at
$\alpha=1/2$. For each fixed $\alpha<1/2$, and every assignment, a
positive proportion on the scale $H/\log x$ of duplicate hits remains
when counting all hits on $S_\alpha$ with multiplicity beyond the first.
To see this, choose four constants
$\alpha<a_1<a_2<b_1<b_2<1/2$. Count pairs with
$x^{a_1}<p\le x^{a_2}$, $x^{b_1}<q\le x^{b_2}$ for which $n$ hits
both rows and has no hit at any prime less than $p$. No assignment
can count a position more than $m(n)-1$ times this way: the first
prime is determined and only the second can vary.
The two classes pick one progression modulo $pq$. Its indices form
an interval of length $H/(pq)+O(1)$, sieved by primes below $p$.
The sieve parameter exceeds $2$ by a fixed amount since
$3a_2+b_2<2$. The lower linear sieve thus gives at least
$cH/(pq\log x)$ positions, with $c>0$ independent of the choices.
The reciprocal sums over both bands tend to positive constants.
This proves the asserted lower bound even for assignments which
leave holes. For each fixed $\alpha>1/2$ a nearly complete cover
instead has $o(H/\log x)$ repeats as proved above.

\paragraph{A test using the additive large sieve.}
Here is one quadratic test of the global profile. Fix $a\in(1/2,1)$
and give each survivor through $x^a$ weight $f(t)$, where
$t=\log p/\log x$ for its first covering prime and $f$ is a bounded
step function with fixed breakpoints. Put
\[
 w(t)=\frac2{t(2-t)},\qquad
 M=\int_a^1 w(t)f(t)\,dt,\qquad B=\int_a^1 tw(t)f(t)\,dt .
\]
Take $Q=\min(x,\sqrt H)/\log x$. Every squarefree $q\le Q$ contains
at most one prime above $x^a$. Write $A$ for the simultaneous CRT
representative of the chosen classes. The Ramanujan sum $c_d(n-A)$
is the constant $\mu(d)$ on the initial survivors whenever $d$ is
squarefree and all its primes are at most $x^a$. For $q=dp$ with a
larger prime, the additional factor is $p\,1_{n\equiv A\bmod p}-1$.
Project the additive Fourier coefficients at each squarefree modulus
onto this one direction by Cauchy--Schwarz. The large sieve upper
bound is $(H+O(Q^2))$ times the sum of the squared weights.
For the lower estimate the row deficits have negligible total mass;
changing weights on multiply hit rows costs at most
$2\|f\|_\infty\sum_{m(n)\ge2}m(n)=o(H/\log x)$.
The uniform row envelopes bound each normalized row by a constant.
Thus the weighted first moments of the row errors also suffice after
squaring. The elementary asymptotic
$\sum_{d\le T}\mu^2(d)/\varphi(d)\sim\log T$ and partial summation
give the following necessary inequality in the limit:
\[
 \int_a^1 wf^2 \ \ge\
 (2-a+\log a)M^2+
 \int_a^1\frac{1-t}{t}\left(\frac{2f(t)}{2-t}-M\right)^2dt .
\]
Indeed the coefficient of $M^2$ from smooth moduli is
$1-\int_a^1(1-t)\,dt/t=2-a+\log a$, since the other moduli each
have a unique large prime and $Q/p<x^a$.
Subtracting the right side from the left gives
\[
   T+(M-B)^2-B^2,\qquad T=\int_a^1\frac{2f(t)^2}{(2-t)^2}\,dt .
\]
Ordinary integral Cauchy--Schwarz gives $B^2\le2(1-a)T\le T$.
Hence these particular large-sieve projections are compatible with
the required profile for every such weight. A contradiction by this
route would need information beyond these projections.

\section*{Coherent assignments and uniform upper bounds}
The upper-bound question admits a deterministic reformulation for
infinite assignments. Give $\mathbb N^{\mathbb N}$ its product topology,
where each coordinate is discrete. A basic open set prescribes finitely
many coordinates. For any cutoff sequence $u_n$ and integer scale $h_n$,
let $G$ consist of the assignments $a$ for which there exist $M,N$ such
that, for every $n\ge N$, the rows of $a$ up to $u_n$ fail to cover
$[1,Mh_n+1]$. Both $M$ and $N$ may depend on $a$. If $G$ is nonmeagre,
there exist fixed $M,N,D$ such that
\[
                      Y(u_n)\le M h_n+D\qquad(n\ge N).
\]
Indeed, for fixed $M,N$ the indicated set is closed, since a cover at
any given $n$ depends on only finitely many coordinates. The countable
union over $M,N$ can be nonmeagre only if one closed set contains a
basic open set. Fixing all rows through some $K$ refines that open set.
At every $n$, the CRT normalization places a cover of length at least
$Y(u_n)-(\prod_{p\le K}p-1)$ inside it. The asserted uniform estimate
follows with $D=\prod_{p\le K}p$. In particular a sequence of upper
bounds proved with assignment-dependent constants on a nonmeagre set
suffices for a global asymptotic upper estimate whenever $h_n\ge1$
eventually.

Conversely, for any $e_n\to\infty$, the assignments reaching
$Y(u_n)-e_n$ infinitely often form a residual set. For each $N$ the
set reaching this threshold at some $n\ge N$ is open. In a prescribed
basic open set the CRT cost is fixed, and some $n\ge N$ has $e_n$ larger
than that cost. This proves density. Intersect these dense open sets
over $N$. Thus the same infinite choice of residues can recur
arbitrarily close to maximality along an infinite subsequence.
An estimate holding with probability one can still leave a meagre
set of good assignments, so the category hypothesis needs its own
argument.
Lean 4.33.1 verifies both \texttt{ProfilsCoherents687.upper\_of\_nonmeagre}
and \texttt{ProfilsCoherents687.near\_optimal\_residual}.

There is also a simultaneous version. For any strictly increasing
sequence of cutoffs $u_0,u_1,\ldots$, one assignment attains $Y(u_0)$
at the first cutoff and covers at least
\[
       Y(u_{j+1})-\Big(\prod_{p\le u_j}p-1\Big)
\]
positions at every subsequent cutoff (with a negative length read as
zero). At step $j+1$ prescribe the preceding assignment through $u_j$.
Each prime row then stabilizes once the cutoff reaches it. Taking its
stabilized value preserves all the earlier covers. For example, along
$u_0=2$ and $u_{j+1}=16^{u_j}$, the loss is at most $\sqrt{u_{j+1}}$ at
every stage. The kernel verifies \texttt{ProfilsCoherents687.simultaneous}.

A long cover also imposes a finite simultaneous condition. Prescribe
the classes below $w$ and let $D=\prod_{p\le w}p$. If
$Y(x)\ge H+TD+(D-1)$, a single assignment with that prefix covers
every length-$H$ window starting at offsets $0,D,\ldots,TD$.
First align a maximal assignment with the prescription by a shift of
at most $D-1$. For the window at $jD$ the classes at all primes shift
by $-jD$, preserving the prescribed rows. In particular
at $z=k^m$, $x=k^{m+1}$ and $H=z^2$, all these translated profiles
must have at most $(24/m)H/\log z$ survivors through $z$ for large $k$.
To exclude the long cover it suffices to prove one translate exceeds
this threshold in each such progression for a single chosen prefix.
Lean verifies \texttt{TranslationsCarrees687.obstruction}.

\section*{Two limits of weighted interval estimates}
Suppose small-prime classes have already been fixed. Put nonnegative
weights on the remaining positions, with total $W>0$. For a remaining
prime $p$, at least one class has weight at least $W/p$. Therefore a
certificate in which the sum of individual prime capacities is strictly
less than $W$ requires $\sum_p1/p<1$. With a fixed small-prime cutoff,
this condition eventually fails by divergence of the prime reciprocal sum.

If the remaining primes are grouped in blocks $P_j$, put
$v_j=\prod_{p\in P_j}(1-1/p)$. Independent uniform choices inside one
block remove expected weight $W(1-v_j)$, so its maximum removal has
at least that size. A certificate obtained by summing the block maxima
therefore requires $\sum_j(1-v_j)<1$. In the regime
$V=\prod_jv_j\le e^{-2}$ this forces one block $i$ to satisfy
\[
 v_i<\tfrac12,\qquad \sum_{j\ne i}(1-v_j)<2V.
\]
If all $v_j\ge1/2$, the inequality $e^{-2(1-v_j)}\le v_j$ gives
$V>e^{-2}$, a contradiction. Choose $i$ with $v_i<1/2$. The sum of
the other deficits is less than $v_i$, so
$\prod_{j\ne i}v_j\ge1-\sum_{j\ne i}(1-v_j)>1/2$. It follows that
$v_i<2V$. This restricts a particular method of proving upper bounds.

\section*{A consequence for prime increments}
Write $J(p)=Y(p)-Y(p-1)$. Since only a prime cutoff introduces a new
class, $Y(x)-Y(z)=\sum_{z<p\le x}J(p)$. The classical Rankin lower
bound gives $Y(x)\gg x\log x\log_3x/(\log_2x)^2$, where
$\log_j$ is the $j$-fold iterated logarithm. Using
$\pi(x)\ll x/\log x$ in the sum shows that, for every $C>0$, there
are arbitrarily large primes $p$ for which
\[
 J(p)>C\left(\frac{\log p}{\log\log p}\right)^2.
\]
Indeed, for any fixed $z$, eventually $2Y(z)\le Y(x)$. Some prime
between $z$ and $x$ then has jump at least $Y(x)/(2\pi(x))$.
The factor $\log_3x$ tends to infinity and $\log t/\log\log t$
is increasing for all sufficiently large $t$. This corollary gives an
unbounded sequence of jumps. Determining the order of their sum remains
the original question.

The Lean declarations are in \texttt{Distribution687}, \texttt{LimitePonderation687},
\texttt{GroupesPonderes687} and \texttt{SautsRankin687}. The Rankin input is
\texttt{BorneRankin687.rankin\_lower}. The underlying problem is recorded
at \url{https://www.erdosproblems.com/687}.
\end{document}
